\section{Control theta del retículo de Leech: identidad \texorpdfstring{\(\Theta_{\Lambda}=E_4^3-720\Delta\)}{Theta Lambda = E4^3 - 720 Delta} y separación entre coincidencia aritmética y morfismo estructural}
\label{sec:control-theta-leech}

Desde \(\Lambda_{24}\), la cadena
\[
\Lambda_{24}
\longrightarrow V^\natural
\longrightarrow\mathbb M
\longrightarrow\text{Moonshine}
\]
es matemática clásica: álgebra de operadores de vértice, grupo Monstruo y
funciones McKay--Thompson. La contribución verificada aquí está en la entrada
finita, la selección del borde y el \(3\)-vecino sin raíces, no en atribuirse de
nuevo esos teoremas.

Tres identidades numéricas acompañan esta cadena:
\[
\Theta_{\Lambda}=E_4^3-720\Delta,
\]
\[
196884=196560+324,
\]
\[
196884=54(3645+1)=54(5\cdot729+1).
\]
Son identidades aritméticas exactas, pero sin un morfismo o una descomposición
de caracteres no constituyen por sí mismas señales matemáticas de
compatibilidad. Ninguna sustituye la construcción de \(V^\natural\) ni
proporciona una acción directa del Monstruo sobre palabras decimales HMT.
